\section{Evaluation: la Measured Capability depende de la Elicitation}

Aunque la evaluación de seguridad y de capacidades parece consistir en ponerle preguntas al modelo, en realidad también abarca el prompt, el scaffold, las tools, el time budget, el environment, el judge y la incertidumbre estadística. La \term{elicitation} (activación de capacidades) consiste en usar prompting adecuado, tool access, search o fine-tuning para que la evaluación se acerque lo más posible a la capacidad que el modelo puede alcanzar. Si la elicitation es demasiado débil, subestima el riesgo; si es demasiado fuerte, puede alejarse de las condiciones reales de un attacker o un user.

\subsection{Task, elicitation, environment y judge}

La evaluación define primero el construct: si mide knowledge, planning, tool use, cyber/CBRN uplift o harmful propensity. Después fija el elicitation budget y los permisos. El environment debe simular el feedback y los constraints reales. El judge puede estar formado por expertos, por un grupo de personas o por modelos, y requiere rubric, inter-rater agreement y calibration.

\lecturefigure{11-evaluation-stack.png}{La puntuación de la evaluación de un modelo depende en conjunto de la task, la elicitation, el environment, el judge y la confidence.}{Anthropic evaluation research; subtítulos locales 27:30--30:40; reconstrucción conceptual.}

\begin{knowledgebox}{Lectura de la figura: una puntuación baja puede indicar poca capacidad o que no se activó}
El fracaso con un prompt fijo no demuestra que el model carezca de la capacidad; el éxito cuando se permiten tools ilimitadas y búsqueda humana tampoco significa que un usuario común pueda reproducirlo. El evaluation report debe indicar model/version, system prompt, tools, samples, search effort, judge, scoring y uncertainty, y distinguir entre observed performance y estimated upper bound.
\end{knowledgebox}

\begin{importantbox}{El intervalo de confianza es más honesto que una puntuación puntual}
Si en $n$ muestras independientes hay $k$ éxitos, la tasa de éxito empírica es
\[
\hat p=\frac{k}{n}.
\]
Aquí $\hat p$ es una estimación con una muestra finita, y la capacidad real $p$ sigue teniendo un intervalo estadístico. Cuando un gate de alto riesgo depende de una failure rate muy baja, unas decenas de muestras son del todo insuficientes. También hay que atender las tareas no independientes, la prompt selection y el evaluator bias; no basta con reportar un porcentaje entero.
\end{importantbox}

\begin{warningbox}{Las evaluaciones de alto riesgo no deben divulgar detalles operativos}
Los materiales públicos pueden describir las categorías de riesgo, el marco de evaluación, el access control, el red-team process y los aggregate findings, pero no deben publicar procedimientos, materiales o métodos de evasión concretos que reduzcan de forma significativa la barrera para conductas dañinas. El objetivo de la transparency es la rendición de cuentas, no convertir el evaluation set en un tutorial de capacidades.
\end{warningbox}

\subsection{Resumen de la sección}

La evaluation es un measurement system, no un banco de preguntas. La elicitation, el environment y el judge determinan qué significa la puntuación; las decisiones de alto riesgo requieren uncertainty, external review y evidencia versionada.
